\section{Conclusion}\label{sec:conclusion}

The emerging-market shortfall in patent survival reflects what gets patented. At equal observable quality and grant
lag, firms in China, Korea and Taiwan keep their patents at least as long as comparable firms in the United States,
Japan, Germany, the United Kingdom and France, and value them about as much relative to foreign owners. On average,
credit shocks and fee cuts barely move their renewal decisions, and they keep the home patents of their inventions
longer than the foreign ones. For firms, the binding constraint is the invention, not the inventor.

Two qualifications matter for policy, and the counterfactuals put a rough price on both. Owner capacity counts at the
margin: patents that small Chinese firms sell to unrelated buyers survive much longer, and moving such patents to
owners able to use them is worth up to \$130--280 million per cohort in China. Policies that deepen markets for
technology, such as registries, licensing intermediaries and secure transfer of title, may therefore add value, while
fee relief, worth almost nothing on the same scale, is unlikely to change which patents are kept. And Chinese
universities are a clear weak spot: their patents are dropped sooner and are worth less than half as much as
comparable foreign ones, consistent with patenting driven by evaluation and subsidy targets rather than use. Yet once
sold to firms, the same patents survive far longer, so part of the loss lies in who holds them: technology-transfer
offices and licensing to firms may recover much of the value, worth \$180--460 million per cohort. Russia, where
domestic patents fall well behind, deserves a closer look.

Three limitations stand out. Only five emerging economies record lapses well enough to study, and two of them, Korea
and Taiwan, are now high-income. The firm results rest on citations to the whole patent family; with citations from
the same office only, emerging-market firms are somewhat behind. And renewal values measure private value to the
owner, not the social value of the knowledge a patent spreads.
